\appendixsection{程序文件说明}\label{sec:appendix-files}

本实验提交的主要程序与结果文件如下：

\begin{itemize}
  \item \texttt{pathfinder.py}：统一实现 BFS、A*、改进 A*、Q-Learning 与人工势场法，并输出静态结果图和 GIF 动态展示；
  \item \texttt{Dijikstra.py}：实现 Dijkstra 算法，并与 BFS、A* 进行单独对比；
  \item \texttt{basic\_map\_comparison.png}、\texttt{challenge\_30x30\_comparison.png}：多算法结果汇总图；
  \item \texttt{challenge\_30x30\_q\_learning\_curve.png}：Q-Learning 训练回报曲线；
  \item \texttt{basic\_map\_bfs.gif}、\texttt{challenge\_30x30\_astar.gif} 等：搜索过程动态展示文件。
\end{itemize}

\appendixsubsection{程序运行方式}

本文所有程序均可在项目根目录下通过以下命令运行：

\begin{verbatim}
python3 pathfinder.py
python3 Dijikstra.py
\end{verbatim}

其中，\texttt{pathfinder.py} 将生成基础场景与复杂场景的对比图、Q-Learning 训练曲线以及 GIF 动画；\texttt{Dijikstra.py} 将生成 Dijkstra 与 BFS、A* 的对比图并输出局限性分析。
